\section{Native IPM Versus Parity IPM}
\label{sec:native-vs-parity-ipm}

The two AC OPF nonlinear routes differ mainly in variable coverage and in how
constraints are represented.

\begin{longtable}{L{0.24\linewidth}L{0.31\linewidth}L{0.33\linewidth}}
\toprule
Feature & Native reduced-space path & Parity full-space path \\
\midrule
\endhead
Primary files & \texttt{ac\_opf.cpp} & \texttt{parity\_formulation.cpp}, \texttt{parity\_ipm.cpp} \\
AC bus variables & non-slack angles and PQ magnitudes & full AC angle and magnitude blocks with bounds \\
Generator dispatch & \(P_g,Q_g\) variables & \(P_g,Q_g\) variables \\
VSC converters & explicit \(P^{ac},Q^{ac},P^{dc}\) variables and balance rows & explicit \(P^{ac},Q^{ac},P^{dc}\) variables and balance rows \\
DC bus voltages & variables, with tight bounds for \texttt{DC\_V} buses & variables with DC balance rows and DC branch limits \\
Load shedding & not modeled & \(\Delta P_d,\Delta Q_d\) variables with VOLL penalty \\
Curtailable renewable / PV & fixed injections & packed decision vectors \texttt{pren\_mw}, \texttt{qren\_mvar} with identity maps \\
Storage & fixed injection & AC and DC storage active dispatch packed into \texttt{pstor\_mw}; AC storage also has \texttt{qstor\_mvar} \\
DCDC converters & not modeled as decision variables & input-side \texttt{pdcdc\_mw} variables folded into DC bus balance physics \\
Flexible loads & projected/fixed demand & \texttt{pflex\_mw} variables with identity maps \\
Failure fallback & fast dispatch plus AC PF if allowed & native path if allowed, then native fallback behavior \\
\bottomrule
\end{longtable}

The parity result stores enhanced component outputs in packed arrays.  The maps
\texttt{ren\_map}, \texttt{stor\_map}, \texttt{dcdc\_map}, and
\texttt{flex\_map} must be used to join those arrays back to the original
component collections.  For storage, \texttt{source\_type=0} denotes AC storage
and \texttt{source\_type=1} denotes DC storage.

One implementation quirk is important for code review: \texttt{ConstraintIndex}
still allocates \texttt{n\_dcdc\_bal} and \texttt{i\_dcdc\_bal}, and
\texttt{build\_problem} includes those rows in \texttt{n\_eq\_total}.  The
current equality evaluator leaves the dedicated DCDC rows unpopulated; DCDC
physics is represented instead through the DC bus balance rows.
